\chapter{Método Proposto}
\label{cap:metodologia}

Metodologia é um termo amplo que descreve um conjunto de procedimentos intelectuais e técnicos para atingir os objetivos.

Neste capítulo, é efetuada uma descrição completa e precisa de todos os procedimentos realizados durante a elaboração e desenvolvimento do trabalho. Devem ser descritas as técnicas, abordagens, instrumentos, equipamentos, programas de computação e outros aparelhos adaptados à pesquisa. O relato deve ser escrito detalhadamente e com clareza, de tal forma que a pesquisa possa ser reproduzida por outro pesquisador e levar aos mesmos resultados.

\section{Classificação da pesquisa}
Esta subseção destina-se a descrever os tipos de pesquisa realizados, classificando-os quanto aos fins (subseção~\ref{sec:fins}) e quanto aos meios (subseção~\ref{sec:meios}). Esta é uma informação importante ao leitor e aumenta a confiabilidade no trabalho.

\subsection{Classificação quanto aos fins} \label{sec:fins}

Quanto aos fins, a pesquisa é descritiva. Segundo, a pesquisa descritiva expõe as características de determinada população ou fenômeno, estabelece correlações entre variáveis e define sua natureza. A autora coloca também que a pesquisa não tem o compromisso de explicar os fenômenos que descreve, embora sirva de base para tal explicação. 

\subsection{Classificação quanto aos meios} \label{sec:meios}

Quanto aos meios, a pesquisa é bibliográfica, telematizada e de campo. A pesquisa é bibliográfica e telematizada pela utilização de teses, dissertações, artigos, livros, jornais e sites na internet para desenvolver e suportar os objetivos propostos nesse estudo. 

A pesquisa é de campo pela utilização de instrumentos como entrevistas em profundidade, além da utilização de questionário.

\section{Limitações do método}
Nesta subseção são discutidas as limitações da pesquisa. É importante comentar que uma pesquisa ter limitações não a invalida, pelo contrário, mostra que o autor sabe os pontos limitadores do estudo e, assim, pode gerar conclusões coerentes.

\section{Universo e amostragem (inaplicável para ''estudo de caso'')}

O universo de pesquisa compreende o conjunto total de elementos que possuem as características relevantes para o estudo, como indivíduos, organizações, eventos ou bases de dados. A amostragem consiste na seleção de uma parte desse universo, de forma criteriosa, com o objetivo de representar adequadamente o todo. A definição do universo e do método de amostragem é essencial para assegurar a validade dos resultados, devendo considerar critérios como acessibilidade, relevância e representatividade dos elementos selecionados.

\section{Coleta e tratamento dos dados [ou das Informações]}

A coleta de dados refere-se aos procedimentos adotados para obtenção das informações necessárias ao desenvolvimento da pesquisa, podendo envolver, por exemplo, análise documental, entrevistas, observação ou uso de bases de dados secundárias. Após a coleta, realiza-se o tratamento dos dados, que consiste na organização, seleção e análise das informações obtidas, com o uso de técnicas adequadas ao tipo de pesquisa. Esse processo visa garantir a consistência, confiabilidade e clareza dos resultados, possibilitando a interpretação adequada e o alcance dos objetivos propostos.

\section{Técnicas utilizadas}

Descreva as técnicas utilizadas, com auxílio de figuras preferencialmente, para que o leitor entenda os princípios de funcionamento de todos os métodos e possa melhor compreender os experimentos.